\documentclass{article}
\usepackage[left=2cm, right=2cm, top=2cm, bottom=2cm]{geometry}
\usepackage{graphicx}
\usepackage{amsmath}
\usepackage{amssymb}
\usepackage{amsthm}
\usepackage{fancyhdr}
\usepackage{verbatim}
\usepackage{listings}
\usepackage{xcolor}
\usepackage{pdfpages}
\usepackage{cancel}
\usepackage{physics}
\usepackage{esint}
\usepackage{braket}

\lstset{
    language=C++,
    basicstyle=\ttfamily\footnotesize,
    keywordstyle=\color{blue},
    commentstyle=\color{green},
    stringstyle=\color{red},
    numbers=left,
    numberstyle=\tiny\color{gray},
    frame=single,
    breaklines=true,
    captionpos=b,
}


\title{MTH 446: Lecture \# 14}
\author{Cliff Sun}

\newtheorem{theorem}{Theorem}[section]
\newtheorem{lemma}[theorem]{Lemma}
\newtheorem{definition}[theorem]{Definition}
\newtheorem{conjecture}[theorem]{Conjecture}
\newtheorem{proposition}[theorem]{Proposition}
\newtheorem{corollary}[theorem]{Corollary}
\newtheorem{one minute paper}[theorem]{One Minute Paper}

\pagestyle{fancy}
\lhead{\textbf{\thepage}\ \ \nouppercase{\rightmark}}
\chead{MTH 446: Lecture \# 14}
\rhead{Cliff Sun}

\begin{document}

\maketitle

\section*{Lecture Span}
\begin{enumerate}
    \item Properties of complex derivatives
\end{enumerate}

\section*{Complex Differentiability}

Remark, if $f \in H(D)$ and $f = u + iv$, then $u,v$ are infinitely differentiable. 

\begin{definition}
    $u$ is said to be \underline{harmonic conjugate} of $u$ if 
    \begin{align}
        u_x &= v_y \\
        u_y &= -v_x
    \end{align}
    In $D$. We abbreviate harmonic conjugate as h.c. 
\end{definition}

In other words, if $u,v \in H(D)$, then $v$ is a h.c. of $u$ in $D$ iff $f = u + iv \in H(D)$. How to find a harmonic conjugate? 

\subsection*{Example}

Let $u = 2x - x^3 + 3xy^2$. Find $v$, the harmonic conjugate of $u$. We can verify that $u$ is harmonic. Use the C.R. equations to obtain that 

\begin{equation}
    v_y = 2 - 3x^2 +  3y^2 \iff v = \int dy (2 - 3x^2 + 3y^2)
\end{equation}

Then, we obtain 

\begin{equation}
    v = 2y - 3x^2y + y^3 + C(x)
\end{equation}

Then 

\begin{equation}
    v_x = -6xy + C'(x) = -u_y = -6xy
\end{equation}

Therefore, $C'(x) = 0 \implies C(x) = \rm const.$ Then, 

\begin{equation}
    v = 2y - 3x^2y + y^3 + C
\end{equation}

\subsection*{Example}

If $v$ is h.c. of $u$ and $\tilde{v}$ is h.c. of $u$, then $\tilde{v} = v + C$ for some const. $C$.  

\subsection*{Example}

Suppose $u,v$ are h.c.'s of each other, then $u,v$ are constant in $D$. 

\begin{proof}
    This means that 
    \begin{align}
        u_x &= v_y \\
        u_y &= -v_x 
    \end{align}
    And vice versa. Then, we arrive that 
    \begin{equation}
        u_x = -u_x \iff u_x = 0, \in D
    \end{equation}
    We obtain similar equations for $v_x, v_y,u_y$. Therefore, $\nabla u \equiv \rm const.$ and $\nabla v \equiv \rm const.$  
\end{proof}

\begin{theorem}
    \underline{Reflection Principle:} Suppose $f \in H(D)$, and contains segment of the x axis and is symmetric to the axis (lower half in $D$ is reflection of the upper half w.r.t the x axis). Then, 
    \begin{equation}
        \overline{f(z)} = f(\overline{z})
    \end{equation}
    if and only if $f(x)$ is real for every $x$ on that segment. 
\end{theorem}

\subsection*{Example}

Suppose $P_n$ is a n-th degree polynomial with real coefficients. Then, one can prove that 
\begin{equation}
    P_n(\overline{z}) = \overline{P_n(z)}
\end{equation}

\begin{definition}
    The exponential of $z$ is 
    \begin{equation}
        \exp(z) = \exp(x)\exp(iy) \iff \exp(x)\left( \cos(y) + i\sin(y) \right)
    \end{equation}
    This is the exponential function.
\end{definition}

\subsection*{Properties}

\begin{enumerate}
    \item $\exp(z_1 + z_2) = \exp(z_1)\exp(z_2)$
    \item $\forall z \in \mathbb{C}$, $|\exp(z)| = e^x > 0$
    \item $\forall z \in \mathbb{C}$, $e^{-z} = 1/e^{z}$. 
    \item $e^{z_1 - z_1} = e^{z_1}/e^{z_2}$
    \item $\exp(z)$ is a periodic function wiht period $p = 2\pi i$ for all $z \in \mathbb{C}$. Note, the period is a complex number.
    \textit{Proof:} $\exp(x)\exp(y + 2\pi) = e^x(\cos(y + 2\pi) + i\sin(2\pi))$. Therefore, $u,v \in C^{\infty}(\mathbb{C})$. 
\end{enumerate}



\end{document}